% ============================================================
% Benchmarking Multi-Agent LLM Architectures for Home Energy Management:
% Real-World Tariff Validation and Cross-Model Cost-Efficiency Analysis
%
% Target: ArXiv preprint (cs.AI + eess.SY cross-list)
% Format: IEEE two-column, 6-8 pages
% ============================================================

\documentclass[conference]{IEEEtran}
\IEEEoverridecommandlockouts

% ── Packages ────────────────────────────────────────────────
\usepackage{cite}
\usepackage{amsmath,amssymb}
\usepackage{graphicx}
\usepackage{booktabs}
\usepackage{hyperref}
\usepackage{xcolor}
\usepackage{multirow}
\usepackage{array}
\usepackage{url}

\graphicspath{{../analysis/figures/}}

% ============================================================
\begin{document}

\title{Benchmarking Multi-Agent LLM Architectures for Home Energy Management:
Real-World Tariff Validation and Cross-Model Cost-Efficiency Analysis}

\author{
  \IEEEauthorblockN{Besnik Sulmataj}
  \IEEEauthorblockA{
    Independent Researcher \\
    \texttt{besniksulmataj@gmail.com}
  }
}

\maketitle

% ============================================================
\begin{abstract}
Multi-agent large language model (LLM) systems have recently been proposed
for home energy management (HEMS), but existing work evaluates only a single
model against simplified tariff structures.
This paper addresses two open questions:
(1) which LLM backend delivers the best cost-performance tradeoff for
multi-agent HEMS scheduling, and
(2) does the agentic HEMS architecture generalize to real-world US utility
tariff structures beyond simplified simulation?
We extend the open-source agentic-ai-hems architecture~\cite{makroum2025agentic}
by benchmarking four LLMs spanning open-source, mid-tier, and frontier tiers
(Llama 4 Maverick, DeepSeek-V3, GPT-4.1, Claude Sonnet 4.6) against three real US utility
time-of-use structures (PG\&E E-TOU-C, SCE TOU-D-4, ComEd Hourly Pricing)
across three household archetypes and 7-day simulations (108 runs total).
Our results show that three of four tested models achieve $>$20\% energy cost
reduction over an unmanaged baseline, with Claude Sonnet 4.6 reaching
49.3\% reduction at \$0.61/day in API cost, and DeepSeek-V3 delivering
37.8\% reduction at only \$0.04/day.
We derive a deployment decision framework mapping LLM API cost against
monthly energy savings by household size and model tier, providing the
first economic viability assessment for agentic HEMS deployment.
\end{abstract}

\begin{IEEEkeywords}
home energy management, large language models, multi-agent systems,
time-of-use tariffs, electric vehicles, energy optimization, cost-efficiency
\end{IEEEkeywords}

% ============================================================
\section{Introduction}
\label{sec:intro}

Residential energy management is undergoing a fundamental shift.
The proliferation of rooftop solar, home battery storage, and electric vehicles
(EVs) has transformed households from passive energy consumers into active
prosumers capable of dynamic load optimization.
Yet the complexity of coordinating these assets --- across time-varying
electricity tariffs, uncertain solar generation, and EV charging deadlines ---
exceeds what rule-based systems handle well.

Large language model (LLM) agents have emerged as a compelling alternative,
offering natural language reasoning about complex, multi-variable scheduling
problems without explicit programming.
Makroum~\cite{makroum2025agentic} demonstrated that a multi-agent LLM
architecture using the ReAct~\cite{yao2023react} pattern can effectively
coordinate appliance scheduling in a HEMS context.
However, that work evaluated only a single model (Cerebras) against
simplified, simulated tariff structures --- leaving two critical questions
unanswered for practitioners considering real-world deployment.

\textbf{Research question:}
\textit{Which LLM backend delivers the best cost-performance tradeoff for
multi-agent home energy management, and does the agentic HEMS architecture
generalize to real-world US utility tariff structures?}

We make three novel contributions:

\begin{itemize}
  \item \textbf{C1 --- First cross-model HEMS benchmark.}
    We compare Llama 4 Maverick, DeepSeek-V3, GPT-4.1, and Claude Sonnet 4.6
    on identical HEMS scheduling tasks across 108 experimental runs,
    reporting energy savings, oracle gap, latency, and token cost per decision.
    Critically, we distinguish between self-hostable open-source models
    (Llama 4 Maverick) and cloud-dependent models, a distinction with direct
    implications for residential energy data sovereignty.

  \item \textbf{C2 --- Real tariff validation.}
    We evaluate against three real US utility TOU structures from 2024
    tariff filings: PG\&E E-TOU-C (3-tier), SCE TOU-D-4 (2-tier),
    and ComEd Hourly Pricing (real-time), representing the dominant
    residential rate designs in the US.

  \item \textbf{C3 --- Economic break-even analysis.}
    We produce the first deployment decision framework for agentic HEMS,
    answering: at what household size does the LLM API cost become
    economically justified by energy savings?
\end{itemize}

% ============================================================
\section{Related Work}
\label{sec:related}

\subsection{Home Energy Management Systems}
The HEMS literature spans a broad range of optimization paradigms.
Zafar et al.~\cite{zafar2020hems} provide a comprehensive survey of HEMS
architectures, cataloguing rule-based, optimization-based, and learning-based
approaches across smart grid deployments.
Traditional MILP formulations~\cite{pinto2023milp} offer provably optimal
scheduling for battery, EV, and solar assets under known price signals,
but require precise system models and become computationally intractable
as the number of controllable assets grows.
Model predictive control (MPC) addresses the receding-horizon nature of
residential scheduling but similarly depends on accurate load and generation
forecasts that are difficult to obtain in practice.
Reinforcement learning (RL) has emerged as a flexible alternative:
Wei et al.~\cite{wei2023drl} demonstrate deep RL for real-time energy
management in smart homes with photovoltaic and storage assets, while
Kaewdornhan et al.~\cite{kaewdornhan2023marl} extend this to multi-home
prosumer settings with EV charging using multi-agent deep RL.
Despite strong results in simulation, RL approaches require substantial
training data, lack interpretability, and struggle to generalize across
tariff structures without retraining.

\subsection{EV and Battery Integration}
Electric vehicle charging has emerged as the dominant flexible load in
residential HEMS~\cite{buijze2024ev}, representing both the largest
controllable demand and the asset with the most stringent scheduling
constraints (departure deadlines, minimum state-of-charge requirements).
Paulauskas and Kapustin~\cite{paulauskas2024battery} demonstrate that
price-arbitrage dispatch of residential battery storage can generate
meaningful cost savings under real-time pricing, but note that naive
scheduling frequently misses optimal windows due to forecast uncertainty ---
a limitation that LLM-based reasoning may help address through contextual
interpretation of price signals.
Wang et al.~\cite{wang2021tou} model residential demand response under
TOU tariffs with stochastic unit commitment, demonstrating that the
structure of the tariff (number of tiers, peak window duration) has a
stronger influence on achievable savings than the specific optimization
method employed.

\subsection{Multi-Agent Systems for Energy}
Multi-agent system (MAS) architectures offer a natural decomposition
for HEMS, where specialist agents manage individual assets and coordinate
through a shared scheduling objective~\cite{binyamin2022mas}.
This mirrors the architecture adopted in this work, where tariff,
solar, battery, and EV agents each maintain local state and report
to a central orchestrating agent.
The key advantage of MAS framing is modularity: new assets (e.g., heat
pumps, V2G) can be added as additional agents without redesigning
the coordination logic.

\subsection{LLM Agents for Planning and Decision-Making}
Beyond the ReAct framework~\cite{yao2023react}, recent work has explored
closed-loop LLM planning agents that adapt their strategies based on
environmental feedback.
Sun et al.~\cite{sun2023adaplanner} propose AdaPlanner, which enables
LLM agents to revise plans in response to execution outcomes ---
a capability directly relevant to multi-step HEMS scheduling where
earlier decisions (e.g., overnight battery charging) affect the
feasibility of later actions.
LLMs have also been evaluated as zero-shot numerical optimizers:
Brahmachary et al.~\cite{brahmachary2024llmopt} show that frontier
models can apply evolutionary reasoning to optimization problems,
though performance varies substantially across model tiers.
Kevian et al.~\cite{kevian2024llmcontrol} benchmark GPT-4, Claude 3 Opus,
and Gemini on control engineering tasks, finding significant performance
gaps between model tiers on problems requiring quantitative reasoning ---
a finding that directly motivates our cross-model HEMS evaluation.
Makroum~\cite{makroum2025agentic} applied the ReAct pattern specifically
to HEMS, demonstrating that LLM agents can coordinate appliance scheduling
without explicit optimization.
Concurrently, Sun et al.~\cite{sun2025llmev} demonstrate LLM-driven
digital twins for EV charging demand response, further establishing
the viability of language models for grid-interactive building control.

\subsection{Data Sovereignty and Privacy in Smart Homes}
A practical barrier to cloud LLM deployment in residential energy
management is data privacy.
Piasecki~\cite{piasecki2024gdpr} documents expert perspectives on GDPR
compliance in smart home IoT deployments, noting that granular energy
consumption data --- including EV charging patterns and occupancy signals ---
constitutes personal data under European regulation.
This concern is amplified in utility-facing deployments where energy
data may be subject to additional sector-specific sharing restrictions.
Self-hosted open-source models such as Llama 4 Maverick represent a
technical response to this constraint, enabling on-premises inference
that eliminates third-party data exposure entirely.

\subsection{Real-World Tariff Structures}
US residential electricity tariffs have diversified rapidly since 2020.
The DOE VGI Assessment~\cite{doe_vgi2025} identifies time-of-use pricing
as the dominant driver of flexible load incentives, with California (PG\&E,
SCE) and Illinois (ComEd) representing the two largest markets with distinct
TOU designs.
The practical implication is that any HEMS solution targeting the US
market must generalize across at least three structurally different
pricing regimes: multi-tier static TOU (PG\&E), simpler two-tier TOU
(SCE), and real-time hourly pricing (ComEd).
To our knowledge, no prior LLM-based HEMS paper has evaluated against
these real utility rate structures.

% ============================================================
\section{Methodology}
\label{sec:methodology}

\subsection{Agent Architecture}
We extend the multi-agent HEMS architecture of~\cite{makroum2025agentic}
with four specialist agents reporting to a central orchestrator:

\begin{itemize}
  \item \textbf{Tariff Agent} --- reports current and forecast electricity
    prices from the parsed utility tariff schedule.
  \item \textbf{Solar Agent} --- reports current generation and 6-hour
    irradiance forecast from the household profile.
  \item \textbf{Battery Agent} --- reports home battery state-of-charge (SoC),
    usable capacity, and charge headroom.
  \item \textbf{EV Charger Agent} --- reports EV SoC, charging urgency,
    and hours until required departure.
\end{itemize}

The \textbf{Orchestrator Agent} receives reports from all four specialists
each hour and produces a JSON decision:

\begin{verbatim}
{
  "ev_charge_rate_pct": <0-100>,
  "home_battery_action": "charge|discharge|idle",
  "solar_action":        "export|store",
  "reasoning":           "<one sentence>"
}
\end{verbatim}

All models are called with \texttt{temperature=0} for reproducibility.
We use the OpenRouter API as a unified inference layer, enabling identical
prompts across all models.
Models were selected to represent four distinct deployment archetypes
rather than to exhaustively cover available options.
The rapidly evolving model landscape makes exhaustive benchmarking
infeasible; we instead characterize the performance-cost tradeoff
across tiers.

\subsection{Tariff Structures}
\label{subsec:tariffs}

We evaluate against three real 2024 US residential electricity tariffs
representing distinct TOU designs:

\begin{itemize}
  \item \textbf{PG\&E E-TOU-C} (California) --- 3-tier structure:
    super off-peak \$0.248/kWh (12am--6am), off-peak \$0.388/kWh,
    peak \$0.574/kWh (4pm--9pm weekdays). Peak-to-valley ratio: 2.3$\times$.

  \item \textbf{SCE TOU-D-4} (California) --- 2-tier structure:
    off-peak \$0.280/kWh, peak \$0.522/kWh (2pm--8pm weekdays).
    No weekend peak. Peak-to-valley ratio: 1.9$\times$.

  \item \textbf{ComEd Hourly Pricing} (Illinois) --- real-time structure
    with prices varying each hour based on day-ahead market.
    Typical range \$0.028--\$0.155/kWh. Evening peak 6pm--8pm.
\end{itemize}

Tariff structures were sourced from 2024 utility tariff filings.
Energy charges only; fixed and demand charges excluded.

\subsection{Household Archetypes}
\label{subsec:households}

We simulate three representative household archetypes (Table~\ref{tab:archetypes})
with synthetic 7-day hourly load profiles generated using fixed random
seeds (42, 43, 44) for reproducibility.
EV daily energy consumption reflects typical US commute distances
($\sim$0.25~kWh/mile).

\begin{table}[h]
\centering
\caption{Household Archetype Specifications}
\label{tab:archetypes}
\begin{tabular}{lccc}
\toprule
 & Small Suburban & Large Suburban & Apartment \\
\midrule
Base load (avg kW)   & 0.5  & 1.2  & 0.3  \\
EV capacity (kWh)    & 60   & 100  & 40   \\
EV daily use (kWh)   & 12   & 15   & 8    \\
Solar peak (kW)      & 5    & 10   & 0    \\
Home battery (kWh)   & 10   & 20   & 5    \\
\bottomrule
\end{tabular}
\end{table}

\subsection{Experiment Design}

We run a full factorial experiment:
$4~\text{models} \times 3~\text{tariffs} \times 3~\text{households}
\times 3~\text{seeds} = 108~\text{runs}$,
each simulating 7 days at hourly resolution (168 scheduling decisions).
Two deterministic baselines are computed for each tariff-household pair:

\begin{itemize}
  \item \textbf{Unmanaged baseline} --- EV charges at maximum rate (7.2~kW)
    immediately upon connection; no battery dispatch; solar output exported
    as generated.
  \item \textbf{Oracle baseline} --- greedy optimal scheduling with perfect
    price foresight: EV charged during cheapest available hours per session;
    battery dispatched via daily arbitrage (charge bottom quartile,
    discharge top quartile).
\end{itemize}

\subsection{Evaluation Metrics}

\begin{align}
\text{Savings} &= \frac{C_{\text{base}} - C_{\text{agent}}}{C_{\text{base}}} \times 100\% \\
\text{Oracle gap} &= \frac{C_{\text{agent}} - C_{\text{oracle}}}{C_{\text{oracle}}} \times 100\% \\
\text{Break-even} &= \frac{C_{\text{API,monthly}}}{C_{\text{base}} - C_{\text{agent}}} \;\; \text{(months)}
\end{align}

where $C_{\text{base}}$, $C_{\text{agent}}$, $C_{\text{oracle}}$ are
7-day electricity costs for the unmanaged, LLM-managed, and oracle
scenarios respectively, and $C_{\text{API,monthly}}$ is the monthly
LLM API cost based on 168 decisions/day at published OpenRouter rates.

% ============================================================
\section{Results}
\label{sec:results}

% ── 4.1 ──────────────────────────────────────────────────────
\subsection{Energy Cost Reduction}

Table~\ref{tab:model_comparison} reports mean energy cost reduction
across all tariff-household combinations (27 observations per model,
3 seeds each). Three of the four tested models exceed the 20\% savings
threshold, validating the agentic HEMS approach across model tiers.
Claude Sonnet 4.6 achieves the highest mean savings at 49.3\%, followed
by GPT-4.1 (44.8\%), DeepSeek-V3 (37.8\%), and Llama 4 Maverick (17.4\%).

\input{../analysis/tables/table1_model_comparison.tex}

All frontier and mid-tier models exceed the 20\% savings threshold,
supporting our hypothesis that the agentic HEMS architecture generalizes
beyond the single-model evaluation of~\cite{makroum2025agentic}.
Figure~\ref{fig:savings} shows per-model distributions.

\begin{figure}[h]
  \centering
  \includegraphics[width=\columnwidth]{fig1_energy_savings}
  \caption{Energy cost reduction vs. unmanaged baseline (\%).
           Error bars show $\pm$1 std across seeds.
           Dashed line indicates 20\% research target.}
  \label{fig:savings}
\end{figure}

% ── 4.2 ──────────────────────────────────────────────────────
\subsection{Oracle Gap Analysis}

The oracle gap measures how close each model's decisions are to the
theoretically optimal schedule with perfect price foresight.
Claude Sonnet 4.6 achieves the smallest oracle gap (31.8\%), closely
followed by GPT-4.1 (34.3\%), indicating that frontier and mid-tier
models capture the majority of available optimization opportunity.
DeepSeek-V3's 45.9\% gap reflects competent but suboptimal scheduling,
while Llama 4 Maverick's 77.8\% gap --- including a 130.1\% gap on
apartment households --- indicates that smaller self-hosted models
occasionally make energy-increasing decisions on complex multi-asset
configurations.

\begin{figure}[h]
  \centering
  \includegraphics[width=\columnwidth]{fig2_oracle_gap}
  \caption{Oracle gap (\% above optimal) by model. Lower is better.}
  \label{fig:oracle_gap}
\end{figure}

Notably, on the large suburban archetype --- the most asset-rich household
with a 100~kWh EV, 20~kWh battery, and 10~kW solar --- Claude achieves a
6.2\% oracle gap, approaching oracle-level performance. This suggests
that frontier models can effectively reason about multi-asset coordination
when given sufficient context.

% ── 4.3 ──────────────────────────────────────────────────────
\subsection{Tariff Generalization}

Table~\ref{tab:tariff_generalization} reports energy savings broken down
by tariff structure.
The ComEd real-time pricing structure presents the greatest opportunity
for all models due to its high intraday price variance
(range \$0.028--\$0.155/kWh), with Claude reaching 84.6\% savings
and even Llama 4 Maverick achieving 36.6\%.
California TOU tariffs prove substantially harder: on PG\&E E-TOU-C,
Llama 4 Maverick achieves only 3.8\% savings --- barely above the
unmanaged baseline --- while Claude achieves 39.0\%.
The 3-tier structure of PG\&E E-TOU-C (with its narrow super off-peak
window from 12am--6am) requires precise overnight scheduling that
smaller models fail to exploit consistently.

\input{../analysis/tables/table2_tariff_generalization.tex}

\begin{figure}[h]
  \centering
  \includegraphics[width=\columnwidth]{fig4_tariff_heatmap}
  \caption{Tariff generalization heatmap: energy savings (\%) across
           model $\times$ tariff combinations.}
  \label{fig:heatmap}
\end{figure}

% ── 4.4 ──────────────────────────────────────────────────────
\subsection{Cost-Efficiency and Break-Even Analysis}

Figure~\ref{fig:cost_efficiency} plots the cost-efficiency frontier ---
energy savings versus daily API cost for each model.
DeepSeek-V3 dominates the cost-efficiency frontier, delivering 37.8\%
savings at only \$0.04/day in API cost --- a 17$\times$ cost advantage
over GPT-4.1 (\$0.27/day) for a 7 percentage-point reduction in
performance.
Claude Sonnet 4.6 occupies the frontier's performance extreme at
\$0.61/day, justified only for high-value households where the
additional 4.5 percentage points of savings over GPT-4.1 translate
to meaningful dollar amounts.

\begin{figure}[h]
  \centering
  \includegraphics[width=\columnwidth]{fig3_cost_efficiency}
  \caption{Cost-efficiency frontier: energy savings vs. daily API cost.
           Points in the upper-left represent the best deployment value.}
  \label{fig:cost_efficiency}
\end{figure}

All models achieve break-even against energy savings within one month
for suburban households (Table~\ref{tab:model_comparison}).
DeepSeek-V3 breaks even in under one week for all household types
at its \$0.04/day API cost.
Claude Sonnet 4.6 on an apartment household --- the least favorable
combination --- breaks even in approximately three weeks given the
apartment's lower energy volume and absence of solar generation.

\begin{figure}[h]
  \centering
  \includegraphics[width=\columnwidth]{fig5_breakeven}
  \caption{Monthly net savings (USD) after deducting API cost, by model
           and household archetype. All models yield positive net savings
           across all household types.}
  \label{fig:breakeven}
\end{figure}

% ============================================================
\section{Discussion}
\label{sec:discussion}

\subsection{Model Tier Recommendations for Deployment}

Our results suggest a tiered deployment strategy based on household
complexity and energy asset portfolio:

\begin{itemize}
  \item \textbf{Apartment / no solar:} GPT-4.1 offers the best tradeoff
    (39.8\% savings, 54.3\% oracle gap, \$0.27/day). Claude's marginal
    performance advantage (37.5\% savings) does not justify its 2.3$\times$
    higher API cost for households with limited solar and smaller batteries.
    Llama 4 Maverick is not recommended for apartments on complex TOU
    tariffs given its near-zero savings on PG\&E structures.

  \item \textbf{Small suburban / moderate solar:} Claude Sonnet 4.6 edges
    GPT-4.1 (56.0\% vs.\ 51.3\% savings; 29.9\% vs.\ 32.8\% oracle gap).
    The 4.7~percentage-point savings advantage represents meaningful
    monthly dollar savings for households with 5~kW solar and 60~kWh EVs,
    making Claude the recommended choice when API budget permits.

  \item \textbf{Large suburban / high solar + EV:} Claude Sonnet 4.6
    achieves a 6.2\% oracle gap --- near-optimal performance --- with
    54.3\% savings vs.\ GPT-4.1's 43.3\% savings and 15.8\% oracle gap.
    For the largest households where absolute dollar savings are highest,
    the frontier model's superior multi-asset coordination justifies
    its cost premium.
\end{itemize}

\subsection{Data Sovereignty and Self-Hosted Deployment}

Llama 4 Maverick's performance --- 17.4\% mean savings, 3.8\% on PG\&E's
complex tariff --- raises important questions for organizations prioritizing
data sovereignty over performance.
As the only self-hostable model in our benchmark, Llama 4 Maverick enables
deployments where residential energy consumption data cannot be transmitted
to third-party cloud APIs --- a requirement increasingly relevant under
GDPR, CCPA, and utility data-sharing regulations.
For ComEd's simpler real-time pricing structure, Llama's 36.6\% savings
approaches DeepSeek's 63.1\%, suggesting that self-hosted models
are viable for markets with high-variance real-time pricing.
However, organizations serving California TOU markets should be aware
that self-hosted models at current capability levels may not reliably
exceed the 20\% savings threshold.

\subsection{Tariff Complexity as a Model Selector}

A practical finding from our tariff breakdown is that tariff complexity
is a stronger model differentiator than household size.
The gap between best and worst model on ComEd Hourly is 48 percentage
points (84.6\% vs.\ 36.6\%), while on SCE TOU-D-4 the gap narrows
to 13 percentage points (24.6\% vs.\ 11.6\%).
PG\&E's 3-tier structure with its narrow super off-peak window is the
most discriminating test: only Claude and GPT-4.1 consistently exploit it,
while Llama essentially fails.
Utility operators and HEMS vendors can use tariff structure complexity ---
specifically the number of pricing tiers and sharpness of peak windows ---
as a proxy for required model capability when selecting deployment tier.

\subsection{Limitations}

This study has several limitations.
First, our oracle baseline uses a greedy daily scheduling heuristic
rather than a true global optimum, which may underestimate the oracle gap.
Second, synthetic household load profiles may not capture the full
variability of real household consumption.
Third, we evaluate a representative summer tariff schedule; seasonal
variation (winter rates, demand charges) is left to future work.
Fourth, 7-day simulation windows, while sufficient for cross-model
comparison, limit our ability to assess long-horizon adaptation.
Finally, token costs are estimated from published OpenRouter rates;
actual costs may vary with caching, context length, and provider pricing changes.

% ============================================================
\section{Conclusion}
\label{sec:conclusion}

We presented the first cross-model benchmark of multi-agent LLM
architectures for residential energy management, evaluated against
real 2024 US utility tariff structures across 108 experimental runs.
Three of four tested models --- DeepSeek-V3 (37.8\%), GPT-4.1 (44.8\%),
and Claude Sonnet 4.6 (49.3\%) --- achieved $>$20\% energy cost reduction
over an unmanaged baseline, validating the agentic HEMS approach beyond
the single-model evaluation of prior work.

DeepSeek-V3 achieves the best cost-efficiency tradeoff at \$0.04/day
API cost, breaking even against energy savings in under one week for
all household types.
Claude Sonnet 4.6 delivers the highest absolute savings (49.3\%) and
comes within 6.2\% of oracle-optimal performance on large suburban
households, justifying its premium for high-value deployments.
GPT-4.1 offers a compelling middle ground: 44.8\% savings at
\$0.27/day with the lowest latency in the benchmark (1,594~ms/decision),
making it the recommended default for latency-sensitive deployments.

These results suggest that agentic LLM scheduling is a commercially
viable approach for residential HEMS when model selection is treated
as an engineering design decision.
The key deployment variable is not model capability per se, but the
match between tariff complexity and model reasoning ability --- a
consideration currently absent from both HEMS literature and commercial
product design.

Future work will extend this analysis to vehicle-to-grid (V2G) incentive
structures, seasonal tariff variation, and multi-dwelling (MDU) deployments.
Longer simulation windows and real household consumption traces would
further strengthen the empirical basis of the cost-efficiency findings
reported here.

% ============================================================
\section*{Acknowledgments}
The author thanks the open-source community behind
\texttt{agentic-ai-hems}~\cite{makroum2025agentic} for providing the
architectural foundation this work extends.
Experiments were conducted using the OpenRouter API.
Tariff data sourced from 2024 utility tariff filings (PG\&E, SCE, ComEd).
All experiments were run on consumer-grade cloud infrastructure, demonstrating the accessibility of this research methodology.

% ============================================================
\bibliographystyle{IEEEtran}
\begin{thebibliography}{99}

\bibitem{makroum2025agentic}
R.~El Makroum,
``Agentic AI for Home Energy Management Systems,''
\textit{arXiv preprint arXiv:2510.26603}, 2025.

\bibitem{yao2023react}
S.~Yao \textit{et al.},
``ReAct: Synergizing Reasoning and Acting in Language Models,''
\textit{arXiv preprint arXiv:2210.03629}, 2022.

\bibitem{doe_vgi2025}
U.S. Department of Energy,
``Vehicle-Grid Integration Assessment,''
Technical Report, January 2025.

\bibitem{pinto2023milp}
P.~M.~E. Pinto and G.~Os\'{o}rio,
``A mixed-integer linear programming approach for optimising residential
energy management with multiple sustainable energy resources,''
\textit{Energy Convers. Manag.}, vol.~278, Art. no. 116594, 2023,
doi: 10.1016/j.enconman.2022.116594.

\bibitem{wei2023drl}
G.~Wei, M.~Chi, Z.-W. Liu, M.~Ge, C.~Li, and X.~Liu,
``Deep Reinforcement Learning for Real-Time Energy Management in Smart Home,''
\textit{IEEE Syst. J.}, vol.~17, no.~2, pp.~2489--2499, Jun. 2023,
doi: 10.1109/JSYST.2023.3247592.

\bibitem{sun2025llmev}
Y.~Sun, C.~Cui, C.~Zhang, and C.~Gong,
``Dynamic Incentive Strategies for Smart EV Charging Stations:
An LLM-Driven User Digital Twin Approach,''
\textit{arXiv preprint arXiv:2504.01423}, Apr. 2025,
doi: 10.48550/arXiv.2504.01423.

\bibitem{zafar2020hems}
A.~Zafar, S.~Bayhan, and A.~Sanfilippo,
``Home Energy Management System Concepts, Configurations, and Technologies
for the Smart Grid,''
\textit{IEEE Access}, vol.~8, pp.~119271--119286, 2020,
doi: 10.1109/ACCESS.2020.3005244.

\bibitem{sun2023adaplanner}
H.~Sun, X.~Zhu, J.~Zhang, D.~Huang, Q.~Qian, and H.~Zhou,
``AdaPlanner: Adaptive Planning from Feedback with Language Models,''
\textit{arXiv preprint arXiv:2305.16653}, 2023,
doi: 10.48550/arXiv.2305.16653.

\bibitem{wang2021tou}
Y.~Wang, U.~Munawar, and R.~Paranjape,
``Stochastic Optimization for Residential Demand Response With Unit
Commitment and Time of Use,''
\textit{IEEE Trans. Ind. Appl.}, vol.~57, no.~2, pp.~1718--1728, 2021,
doi: 10.1109/TIA.2020.3048643.

\bibitem{buijze2024ev}
M.~Buijze \textit{et al.},
``Enhancing smart charging in electric vehicles by addressing paused
and delayed charging problems,''
\textit{Nature Commun.}, vol.~15, 2024,
doi: 10.1038/s41467-024-48477-w.

\bibitem{paulauskas2024battery}
L.~Paulauskas and N.~Kapustin,
``Battery Scheduling Optimization and Potential Revenue for Residential
Storage Price Arbitrage,''
\textit{Batteries}, vol.~10, no.~7, p.~251, 2024,
doi: 10.3390/batteries10070251.

\bibitem{binyamin2022mas}
S.~S. Binyamin and S.~A.~C. Slama,
``Multi-Agent Systems for Resource Allocation and Scheduling in a Smart Grid,''
\textit{Sensors}, vol.~22, no.~21, p.~8099, 2022,
doi: 10.3390/s22218099.

\bibitem{kevian2024llmcontrol}
D.~Kevian \textit{et al.},
``Capabilities of Large Language Models in Control Engineering:
A Benchmark Study on GPT-4, Claude 3 Opus, and Gemini 1.0 Ultra,''
\textit{arXiv preprint arXiv:2404.03647}, 2024,
doi: 10.48550/arXiv.2404.03647.

\bibitem{piasecki2024gdpr}
S.~Piasecki,
``Expert perspectives on GDPR compliance in the context of smart homes
and vulnerable persons,''
\textit{Inf. Commun. Technol. Law}, vol.~33, no.~1, 2024,
doi: 10.1080/13600834.2023.2231326.

\bibitem{kaewdornhan2023marl}
N.~Kaewdornhan \textit{et al.},
``Real-Time Multi-Home Energy Management with EV Charging Scheduling
Using Multi-Agent Deep Reinforcement Learning Optimization,''
\textit{Energies}, vol.~16, no.~5, p.~2357, 2023,
doi: 10.3390/en16052357.

\bibitem{brahmachary2024llmopt}
S.~Brahmachary \textit{et al.},
``Large Language Model-Based Evolutionary Optimizer: Reasoning with Elitism,''
\textit{arXiv preprint arXiv:2403.02054}, 2024,
doi: 10.48550/arXiv.2403.02054.

\end{thebibliography}

\end{document}
